%%%%%%%%%%%%%%%%%%%%%%%%%%%%%%%%%%%%%%%%%%%%%%%
% Wartungsplan CNC-Drehbank Wabeco CC-D6000 - Jahr (Formular zum Ausdrucken)
% CC-BY-SA 3.0, FAU FabLab
% Analog zum Wartungsplan der CNC-Fräse (fraese-einweisung).
%%%%%%%%%%%%%%%%%%%%%%%%%%%%%%%%%%%%%%%%%%%%%%%
\PassOptionsToPackage{table}{xcolor}
\documentclass[a4paper,landscape,10pt]{scrartcl}
\usepackage{wartungsplan}
\renewcommand{\PlanPfad}{Dateifreigabe/1\_FabLab\_Einweisungen/drehbank-einweisung/Wartungsplan\_Jahr.pdf}

\begin{document}

\Kopf{Wartungsplan CNC-Drehbank Wabeco CC-D6000 -- Jahr}{%
  Jahr: 20\rule{10mm}{0.4pt}
  \qquad Plan begonnen am: \rule{25mm}{0.4pt}}
\small
In die \textbf{farbigen Felder} Datum und Kürzel eintragen, wenn die Aufgabe erledigt ist.
Graue Felder: nicht fällig. Die wöchentlichen Arbeiten stehen auf Seite 2.\\[1mm]
\Legende

\vspace{2.5mm}
% ------------------------------ Plan ---------------------------------
\Spalten{12}
\begin{tikzpicture}[x=1pt,y=1pt]
  % Kopfzeile des Plans: Quartale und Monate
  \pgfmathsetlengthmacro{\xa}{\wG+\wI+\wT}
  \node[anchor=west, font=\small\bfseries] at (1.2mm, 9mm) {Aufgabe};
  \node[anchor=west, font=\footnotesize] at (1.2mm, 3.2mm) {Intervall und Tätigkeit};
  \node[anchor=east, font=\footnotesize] at (\xa-1.2mm, 3.2mm) {Erledigt am (Datum):};
  \foreach \q [count=\k from 0] in {Q1,Q2,Q3,Q4}{
    \node[font=\small\bfseries] at (\xa+3*\k*\wS+1.5\wS, 9mm) {\q};
  }
  \foreach \m [count=\k from 1] in {Jan,Feb,Mär,Apr,Mai,Jun,Jul,Aug,Sep,Okt,Nov,Dez}{
    \node[font=\footnotesize] at (\xa+\k*\wS-0.5\wS, 3.2mm) {\m};
  }
  \draw[black, line width=1.2pt] (0,0) -- (\xa+12\wS,0);
  \setcounter{zeile}{0}
  % Reinigung und Schmierung
  \Aufgabe{iM}{monatlich}{}{xxxxxxxxxxxx}
  \Aufgabe{iM}{monatlich}{}{xxxxxxxxxxxx}
  \Aufgabe{iM}{monatlich}{}{xxxxxxxxxxxx}
  \Aufgabe{iM}{monatlich}{}{xxxxxxxxxxxx}
  \Aufgabe{iQ}{vierteljährlich}{}{x..x..x..x..}
  \Aufgabe{iQ}{vierteljährlich}{}{x..x..x..x..}
  \Aufgabe{iQ}{vierteljährlich}{}{x..x..x..x..}
  \Aufgabe{iQ}{vierteljährlich}{}{x..x..x..x..}
  \Aufgabe{iH}{halbjährlich}{}{x.....x.....}
  \Aufgabe{iJ}{jährlich}{}{x...........}
  \Gruppe{1}{10}{Reinigung /\\ Schmierung}
  \Trenner
  % Bestätigung / Prüfung
  \Aufgabe{iQ}{vierteljährlich}{}{x..x..x..x..}
  \Aufgabe{iH}{halbjährlich}{}{x.....x.....}
  \Aufgabe{iH}{halbjährlich}{}{x.....x.....}
  \Aufgabe{iH}{halbjährlich}{}{x.....x.....}
  \Aufgabe{iJ}{jährlich}{}{x...........}
  \Gruppe{11}{15}{Bestätigung /\\ Prüfung}
  \Trenner
  % Austausch
  \Aufgabe{iJ}{jährlich}{}{x...........}
  \Gruppe{16}{16}{\tiny Tausch}
  \Trenner
  % Sonstiges: frei für weitere Aufgaben
  \Aufgabe{}{}{}{------------}
  \Aufgabe{}{}{}{------------}
  \Gruppe{17}{18}{Sonst.}
  \Rahmen{3,6,9}
\end{tikzpicture}

\newpage
% ------------------------ Wöchentliche Arbeiten ----------------------
{\Large\bfseries Wöchentliche Arbeiten} \hfill \small Wartungsplan CNC-Drehbank Wabeco CC-D6000 -- Jahr\\[1mm]
\normalsize
Auszuführen am \textbf{Donnerstag}. Kalenderwoche durchstreichen, wenn alles erledigt wurde.

\vspace{2mm}
\begin{tikzpicture}[x=1pt,y=1pt]
  \pgfmathsetlengthmacro{\wk}{\linewidth/27}
  \foreach \i in {1,...,53}{
    \pgfmathtruncatemacro{\r}{int((\i-1)/27)}
    \pgfmathtruncatemacro{\s}{\i-1-27*\r}
    \draw[iW!35, fill=iW!35] (\s*\wk,-\r*10mm) rectangle ++(\wk,-10mm);
    \draw[black!60, line width=0.3pt] (\s*\wk,-\r*10mm) rectangle ++(\wk,-10mm);
    \node[font=\footnotesize] at (\s*\wk+0.5*\wk,-\r*10mm-5mm) {\i};
  }
  \draw[black, line width=1.2pt] (0,0) rectangle (27*\wk,-20mm);
\end{tikzpicture}

\vspace{6mm}
\Unregelmaessigkeiten{Wartungsplan CNC-Drehbank Wabeco CC-D6000 -- Jahr}{9}

\end{document}
